\chapter{Conclusions and Outlook} \label{Conclusions}
This thesis has shown that predictive control based on qLPV models is fast enough for real-time wind turbine control and improves power reference tracking compared to established baseline schemes. Furthermore, it was demonstrated that Koopman-based MPC algorithms achieve smooth power reference tracking in wind farm control. Section~\ref{sec:Conclusions} summarizes these results, and Section~\ref{sec:Outlook} discusses resulting directions for future research.
\section{Conclusions}\label{sec:Conclusions}
This dissertation develops qLPV, structured low-order and data-driven control schemes for wind energy conversion systems at the turbine and wind farm levels. The aim is to minimize the Levelized Cost of Energy (LCoE) of increasingly large and slender turbines in dense wind farms by balancing power yield, actuator usage, and mechanical loads. For turbine control, velocity-based qLMPC handles the nonlinearity of the turbine models. For farm control, data-driven MPC based on Koopman operator theory is used, as the high-dimensional ODEs of the discretized flow equations prevent real-time physics-based MPC. All algorithms are real-time capable, with mean computation times well below the typical turbine controller sample time of about 10\,ms.

One contribution of this thesis is the derivation of qLPV wind turbine models with a good fit relative to FAST in the range of 0.1\,Hz to 0.7\,Hz, which is relevant for turbine control. Two models are designed, a 5th order qLPV model with rotor and tower dynamics, and a 9th order model containing additionally blade out-of-plane, generator and slip angle states. The latter fits better in open loop but brings only marginal closed-loop gains. It was nevertheless chosen as the qLMPC plant, as it causes no measurable computational penalty and contains a subset of states close to an established prediction model \cite{Bianchi2007}; for the application to FAST, the prediction model was extended to all its states except the tower side-to-side motion. The velocity-based MIMO qLMPC was benchmarked against a collective pitch and torque SISO baseline, together with Individual Pitch Control (IPC) and aerodynamically modeled Trailing Edge Flap Control (TEFC); all configurations maintain the mean power. Without model mismatch, the qLMPC reduces the standard deviations of the tower fore-aft acceleration and the power tracking error by factors of up to 4.1 and 11.4. On FAST, the adapted designs still reduce the power fluctuations by a quarter to a third, which no other approach achieves, but increase the tower Damage Equivalent Load (DEL) and do not reduce the blade DEL. IPC and the combined IPC-TEFC, referred to as Active Aerodynamic Load Control (AALC), reduce the blade DEL above rated by 21\% to 24\% on the optimization and two verification datasets with gains optimized at a single wind speed, AALC at a considerably lower pitch actuator power. Only the combinations of qLMPC and IPC reduce both the blade DEL and the power fluctuations. These results underline the potential of qLMPC, but also its dependence on an accurate representation of the wind acting on the rotor and of the turbine dynamics.

At the farm level, this thesis introduces a data-driven adaptive Axial Induction Control (AIC) in WFSim, based on Koopman operator theory: the model is identified via Dynamic Mode Decomposition and physics-inspired lifting functions. It estimates the effective wind speeds and tracks power references that a non-adaptive scheme trained on the same, deliberately incomplete dataset cannot reach. Updating the model only when significant discrepancies are detected reduces the tracking error from 456\,kW to 17.2\,kW, with a Variance-Accounted-For of the effective wind speed of nearly 100\% for the upstream turbine and of about 80\% in open loop and 98\% in closed loop for the downstream turbine. Combined Wake Redirection Control (WRC) and AIC was explored in two ways: pre-optimizing the yaw misalignment with the Bastankhah-Gaussian wake model, and a WRC-AIC Koopman MPC optimizing thrust and yaw jointly. With a linear Koopman model lifted by filtered control signals, nonlinear control features and moving means, and using only the measured turbine powers, the joint optimization reduces the tracking error for power references above the greedy power from 118\,kW with baseline AIC and 65.3\,kW with pre-optimized yaw to 27.8\,kW, showing that data-driven models can capture the physical effects to optimize the farm yield.
\section{Outlook} \label{sec:Outlook}
%Future research can follow three directions: bringing the simulation models and the evaluation closer to reality, validating the controllers in field tests, and investigating neural network-based models.

Despite the promising results, several idealized assumptions within the simulation framework must be addressed to transition to real-world application. The individual turbine simulations assume that the wind at the turbine is perfectly known. While wind estimators are routinely used in industry \cite{soltani2013estimation}, their accuracy is inherently sensitive to turbine parameters that inevitably change as the system ages: uncertainty in the power coefficient, for example, biases the estimated rotor-effective wind speed \cite{brandetti2022ill}. Sensors such as spinner anemometers and lidar offer an alternative. The assumption of a perfectly known current wind speed is optimistic, whereas not using any wind preview is pessimistic: with lidar preview, both controllers in \cite{schlipf2012comparison} considerably reduce the tower DEL compared to a baseline feedback controller, whereas the qLMPC in this thesis, using only the current wind speed, increases it on FAST. Future work should therefore explore lidar-based wind preview to enhance MPC schemes for both Collective Pitch Control and IPC, particularly for gust alleviation.
%
Moreover, the application of the qLMPC to FAST showed that its tower load reduction is lost, probably because of the single-point wind speed used for scheduling, while the rotor averages a full wind field, and because of dynamics above about 0.7\,Hz that are not contained in the prediction model. That model mismatch limits the advantage of MPC is also suggested by \cite{schlipf2012comparison}, where a simple feedforward controller combined with collective pitch and tower feedback performs slightly better than a nonlinear MPC based on a reduced model, although its additional tower feedback also contributes. A further source of model mismatch is the actuator model: the pitch and generator torque dynamics are modeled as first-order transfer functions with constant parameters. In practice, these actuator responses degrade over time, so that learning and adapting to changed actuator behavior is a necessary research extension. 
%
The wind farm studies considered a wind speed ramp and constant inflow without turbulence. While the adaptive AIC presented here was shown to be adequate to address this task, the constant forgetting factor could be replaced by the matrix forgetting factor used for Koopman-based MPC in \cite{calderon2023data}. The results presented in~\cite{dittmer2026koopman} extend the WFSim environment to support turbulent inflow using the Kaimal Normal Turbulence Model, demonstrating that while turbulence reduces the potential gains from WRC compared to steady-state predictions, the combined WRC-AIC scheme still outperforms baseline AIC, reducing power tracking error considerably under realistic, time-varying wind conditions. Closed-loop WRC in this work relies on the measured turbine powers only, but load-based wake tracking could provide the wake position as an additional feedback signal. Such load-based estimates were validated against scanning lidar measurements in a field test in \cite{onnen2026field}.
%
The load evaluation under WRC-induced yaw misalignment remains an open problem, which cannot be addressed with WFSim's actuator disk turbine model. This quasi-static turbine model, which represents the turbine as an energy sink without mechanical dynamics apart from a first-order filter on the thrust control signal, is also not capable of imitating the dynamic behavior required to fully evaluate power demand following. Transitioning to environments such as FAST.Farm would allow for the inclusion of dynamic turbine behavior, a comprehensive evaluation of loads during WRC, and the investigation of Dynamic Induction Control alongside AIC and WRC. Furthermore, the framework should be tested on a higher number of turbines under realistic, varying wind directions.
%
The evaluation of the controllers should also become more realistic. Another research opportunity lies in defining a concrete financial advantage for load alleviation against the disadvantage of increased actuator wear. Quantifying this trade-off is essential for minimizing the LCoE with a commercially viable controller, see, e.g., \cite{pusch2024impact}. The actuator power criterion used in this thesis is only a proxy for actuator wear: full-scale tests of blade bearings show that the wear caused by IPC depends strongly on the bearing design, the lubrication, and lubrication runs, and that tests with static loads can be harsher than those with dynamic loads \cite{bartschat2023effect}.

To prove industrial validity, the controllers should ultimately be validated with measurement data. Research wind farms such as WiValdi, operated by the German Aerospace Center (DLR) together with its partners in the Research Alliance Wind Energy, provide the required instrumentation: inflow measurements from meteorological masts and lidar, together with load measurements on instrumented turbines, would allow both the effective wind speed estimates and the load reductions predicted here to be verified experimentally.

Independently of these steps, neural network-based models offer a promising alternative for the data-driven farm control. The physically motivated lifting functions selected in this thesis were shown to yield accurate Koopman models for wind farm control. Subsequent work~\cite{sharan2024autoencoder, sharan2026tcn}, to which the author contributed, has explored neural network-based lifting functions, suggesting that learned lifting functions can outperform manually selected ones. This direction is promising for larger farms and more complex inflow, where overlapping wakes, direction-dependent interactions, and varying transport delays make suitable lifting functions hard to select based on physical insight. As the MPC acts on the lifted model, learned lifting functions are part of the control law. While the input constraints remain satisfied, guaranteeing safe behavior of the predicted states is difficult as long as the behavior of the network outside its training data cannot be bounded.

%The results highlight the potential of physics-based and data-driven predictive multivariable control for the increasingly complex trade-offs in wind energy production.